\documentclass{notebook}
\begin{document}

\section{Chi-square testによるモデル適合度の評価}
モデルから計算される光度曲線が観測データを説明できるか
どうかを評価するため、chi-square testを行った。

観測データを \(d=(y_i^{\mathrm{obs}},\sigma_i)_{i=1}^{N}\)、
モデルから得られる値を \(y_i^{\mathrm{model}}(\theta)\)とする。
$\theta$は自由パラメータであり、今回は\(\theta = (A_{\wind},\beta_{\mathrm{sh}})\)である。

$\chi^2(\theta)$を以下で定義する：
\begin{equation}
  \chi^2(\theta)
  \coloneq
  \sum_{i=1}^{N}
  \ab(
      \frac{y_i^{\mathrm{obs}}-y_i^{\mathrm{model}}(\theta)}{\sigma_i}
  )^2
\end{equation}
\end{document}
